\begin{table}[htbp]
    \centering
    \caption{Representative convergence and runtime results for CG, BJ-PCG, and Classical AS-PCG.}
    \label{tab:representative-results}
    \small
    \setlength{\tabcolsep}{3.5pt}
    \begin{tabular}{rrlrrrr}
        \toprule
        \(N\) & \(p\) & Method & \(N_{\mathrm{it}}\) & \(r_{\mathrm{rel}}\) & \(T_{\mathrm{total}}\) (s) & \(T_{\mathrm{iter}}\) (ms) \\
        \midrule
        \multicolumn{7}{l}{\textit{Single-process reference}} \\
        1024 & 1 & CG & 1898 & \(9.959\times10^{-9}\) & \(10.821\pm0.023\) & 5.701 \\
        1024 & 1 & BJ-PCG & 383 & \(9.329\times10^{-9}\) & \(23.062\pm0.693\) & 60.215 \\
        1024 & 1 & AS-PCG (\(\delta=1\)) & 383 & \(9.329\times10^{-9}\) & \(26.980\pm0.562\) & 70.445 \\
        1024 & 1 & AS-PCG (\(\delta=2\)) & 383 & \(9.329\times10^{-9}\) & \(26.981\pm0.723\) & 70.447 \\
        \addlinespace
        \multicolumn{7}{l}{\textit{Large-problem strong-scaling configuration}} \\
        2048 & 32 & CG & 3826 & \(9.848\times10^{-9}\) & \(3.443\pm0.019\) & 0.900 \\
        2048 & 32 & BJ-PCG & 1102 & \(9.235\times10^{-9}\) & \(10.527\pm1.103\) & 9.552 \\
        2048 & 32 & AS-PCG (\(\delta=1\)) & 1072 & \(9.488\times10^{-9}\) & \(12.650\pm1.782\) & 11.801 \\
        2048 & 32 & AS-PCG (\(\delta=2\)) & 1101 & \(9.992\times10^{-9}\) & \(13.370\pm1.091\) & 12.143 \\
        \addlinespace
        \multicolumn{7}{l}{\textit{Largest weak-scaling configuration}} \\
        2896 & 128 & CG & 5429 & \(9.962\times10^{-9}\) & \(3.173\pm0.465\) & 0.584 \\
        2896 & 128 & BJ-PCG & 2135 & \(9.622\times10^{-9}\) & \(16.077\pm2.557\) & 7.530 \\
        2896 & 128 & AS-PCG (\(\delta=1\)) & 1900 & \(9.502\times10^{-9}\) & \(19.321\pm2.798\) & 10.169 \\
        2896 & 128 & AS-PCG (\(\delta=2\)) & 1909 & \(9.949\times10^{-9}\) & \(22.352\pm2.942\) & 11.709 \\
        \bottomrule
    \end{tabular}
    \par\smallskip
    \begin{minipage}{0.96\textwidth}
        \footnotesize
        Reported runtimes are means \(\pm\) standard deviations over five runs. The per-iteration time is calculated as \(T_{\mathrm{iter}}=1000T_{\mathrm{total}}/N_{\mathrm{it}}\).
    \end{minipage}
\end{table}
